\id{ҒТАМР 65.63.33}{}

\begin{header}
\swa{В.~Kalemshariv, R.~Askar, А.К.~Mustafayeva, М.Т.~Mursalykova, А.М.~Shulenova}{TECHNOLOGY FOR PRODUCING KURT FROM CAMEL MILK}

В.~Kalemshariv,
R.~Askar,
А.К.~Mustafayeva\envelope,
М.Т.~Mursalykova,
А.М.~Shulenova
\end{header}

\begin{affil}
NCJSC «S.Seifullin Kazakh Agro Technical Research University», Kazakhstan

\corrauthor{Corresponding-author: ayaulym.mustafa@mail.ru}
\end{affil}

The article is devoted to the preparation of kurut, a national dairy
product made from camel milk. A technology has been developed for using
plant raw materials to enrich national dairy products with vitamins.
Camel milk from Southern Kazakhstan and rose hips from Eastern
Kazakhstan were selected as research objects, with rose hips used as
plant raw materials added to the finished curd mass in dried, powdered
form. This is because the vitamin content in dried fruits and berries is
preserved for a long period of time. Rose hips are rich in vitamins and
are especially valuable as a primary source of ascorbic acid, containing
ten times more than apples. Ascorbic acid plays a key biochemical and
physiological role in the human body; it contributes to the normal
development of connective tissues, stress resistance and normal immune
function, and protects against viral and bacterial infections. Three
samples of cow' s milk and three samples of
camel' s milk were taken for the study, as well as a
control sample and samples of cottage cheese enriched with 3\% and 5\%
rosehip powder, and an analysis of the vitamin C enrichment level was
carried out. The physicochemical, organoleptic and energy indicators of
cottage cheese with added dry rosehip powder were calculated. The sample
with 5\% rosehip powder, which showed the best results, was recognised
as optimal. The article is devoted to the preparation of kurut, a
national dairy product made from camel milk. A technology has been
developed for using plant raw materials to enrich national dairy
products with vitamins. Camel milk from Southern Kazakhstan and rose
hips from Eastern Kazakhstan were selected as research objects, with
rose hips used as plant raw materials added to the finished curd mass in
dried form. The vitamin C content in 100 g of the product was determined
to be 11,754 ± 0,24\%, protein-60,1, fat-3,1, carbohydrates-18,2, energy
value-341,1 kcal.

{\bfseries Keywords:} cow' s milk, camel milk, national
product, kurt, vitamin C, rose hips.

\begin{header}
ТЕХНОЛОГИЯ ПРОИЗВОДСТВА КУРТА ИЗ ВЕРБЛЮЖЬЕГО МОЛОКА

Б.~Калeмшарив,
Р.~Асқар,
А.К.~Мустафаeва\envelope,
М.Т.~Мурсалыкова,
А.М.~Шуленова
\end{header}

\begin{affil}
НАО «Казахский агротехнический исследовательский университет им. С.Сейфуллина», Казахстан,

е-mail: ayaulym.mustafa@mail.ru
\end{affil}

Статья посвящена приготовлению национального молочного продукта-курута
из верблюжьего молока. Разработана технология использования
растительного сырья для обогащения национальных молочных продуктов
витаминами. В качестве объектов исследования были выбраны верблюжье
молоко из Южного Казахстана и шиповник из Восточного Казахстана, причем
шиповник, используемый в качестве растительного сырья, добавлялся в
готовую творожную массу в сушеном, порошкообразном виде. Это связано с
тем, что содержание витаминов в сушеных фруктах и ягодах сохраняется в
течение длительного периода времени. Шиповник богат витаминами, особенно
ценен как основной источник аскорбиновой кислоты, содержащей в десять
раз больше, чем яблоко. Аскорбиновая кислота играет ключевую
биохимическую и физиологическую роль в организме человека; она
способствует нормальному развитию соединительных тканей,
стрессоустойчивости и нормальной иммунной функции организма, а также
защищает от вирусных и бактериальных инфекций. Для исследования были
взяты три образца коровьего молока и три образца верблюжьего молока, а
также контрольный образец и образцы творога, обогащенного 3\% и 5\%
порошком шиповника, и был проведен анализ уровня обогащения витамином С.
Были рассчитаны физико-химические, органолептические и энергетические
показатели творога с добавлением сухого порошка шиповника. Образец с 5\%
порошком шиповника, который показал лучшие результаты, был признан
оптимальным. Содержание витамина С в 100 г продукта было определено как
11,754 ± 0,24\%, белок-60,1, жир-3,1, углевод-18,2, энергетическая
ценность-341,1 ккал.

{\bfseries Ключевые слова:} коровье молоко, верблюжье молоко, национальный
продукт, курт, витамин С, шиповник.

\begin{header}
ТҮЙЕ СҮТІНЕН ҚҰРТ ӨНІМІН ДАЙЫНДАУ ТЕХНОЛОГИЯСЫ

Б.~Калeмшарив,
Р.~Асқар,
А.К.~Мустафаeва\envelope,
М.Т.~Мурсалыкова,
А.М.~Шуленова
\end{header}

\begin{affil}
«С.Сейфуллин атындағы Қазақ агротехникалық зерттеу университеті» КеАҚ, Қазақстан,

е-mail: ayaulym.mustafa@mail.ru
\end{affil}

Мақала түйе сүтінен ұлттық сүт өнімін-құртты дайындауға арналған.Ұлттық
сүт өнімдерін дәрумендермен байыту мақсатында өсімдік тектес шикізаттар
қолданудың технологиясы әзірленді. Зерттеу нысандары ретінде Оңтүстік
Қазақстанда өсірілген түйе сүті және өсімдік шикізаты ретінде Шығыс
Қазақстанда өсірілген итмұрын таңдап алынып, дайын болғаны сүзбе
массасына құрғақ және ұнтақталған күйінде қосылды. Өйткені құрғақ күйде
кептірілген жеміс-жидектердегі дәрумендердің мөлшері ұзақ уақыт
аралығында сақталады. Итмұрын дәрумендерге бай, соның ішінде ең бірінші
аскорбин қышқылының көзі ретінде құндылығы өте жоғары, алма жемісімен
салыстырғанда 10 есе көп болады. Аскорбин қышқылы адам ағзасында негізгі
биохимиялық және физиологиялық рөл атқарады, ол дәнекер тіндердің
қалыпты дамуына, стресске төзімділікке, дененің қалыпты иммундық
жағдайын қалыптастыруға, вирустық және бактериялық инфекциядан қорғауға
ықпал етеді. Зерттеу жүргізу үшін сиыр сүтінің үш үлгісі, түйе сүтінің
үш үлгісі алынып, бақылау үлгісі және 3\%, 5\% итмұрын ұнтағы қосылған
құрт үлгілерінің С дәруменімен байытылу дәрежесі зерттелді. Зерттеу
нәтижесінде құрғақ ұнтақталған итмұрын қосылған құрттың физика-химиялық
көрсеткіштері, органолептикалық көрсеткіштері және энергетикалық
құндылығы есептелді. Жоғары нәтиже көрсеткен 5\% итмұрын ұнтағы қосылған
үлгі оңтайлы деп тынылды.100 гр өнімдегі С дәруменнің мөлшері
11,754±0,24\%, ақуыз-60,1, май-3,1, көмірсу-18,2, энергетикалық
құндылығы-341,1 ккал болатындығы анықталды.

{\bfseries Түйін сөздер}: сиыр сүті, түйе сүті, ұлттық өнім, құрт, С
дәрумені, итмұрын.

\begin{multicols}{2}
{\bfseries Introduction.} Currently, the global dairy market is actively
developing, and the assortment of manufactured products is expanding day
by day. However, due to the economic situation and changes in the
international environment, new approaches and methods that ensure the
production of high-quality products in the dairy production and
processing sector are still being considered. One quarter of Kazakhstan
is characterized as flat steppe, another quarter as foothill regions,
while half of the country is considered desert and semi-desert areas
where camel breeding plays a special role. As of the end of 2024, the
number of camels in Kazakhstan exceeded 159,111 thousand, and there is a
need to regulate industrial processing of camel milk into
export-oriented products. For Kazakhstan, in order to further develop
the sector, one of the main tasks at present is the industrial
processing and production of camel milk, which is currently considered
at a low level. The production of camel milk, as well as the preparation
and consumption of products such as cheese and yogurt from it, has been
growing significantly in recent years {[}1{]}.

Camel milk differs from the milk of other agricultural animals in its
chemical composition, nutritional and medicinal properties. Camel milk
contains three times more vitamin C and ten times more iron than cow's
milk, as well as unsaturated fatty acids, B vitamins, and mineral
substances. The average chemical composition of camel milk dry matter is
8.2\%, fats-3.5\%, proteins-2.9\% (including casein-2.5\%),
carbohydrates-4.7\%. Camel milk is also a rich source of calcium,
phosphorus, and fat-soluble vitamins {[}2{]}.

Known for its immune-forming, anti-inflammatory, apoptotic, and
antidiabetic properties, camel milk is considered a naturally beneficial
food. Due to its high content of β-casein and various secreted
antibodies, it is easily digestible and, compared to cow's milk, is
capable of resisting bacteria and viruses. The β-casein in camel milk
does not cause allergies and is well absorbed, as it is sensitive to
gastrointestinal hydrolysis; therefore, due to the high level of
β-casein, camel milk is beneficial to human health {[}3{]}.

Kazakh national dishes represent a traditional diet formed over many
centuries. The daily diet of nomads is very beneficial to health and is
distinguished by hospitality and generosity unique to the Kazakh people.
Based on the main ingredients in their composition, Kazakh national
dishes are divided into four groups: flour-based dishes, cereal-based
dishes, meat dishes, and dairy dishes.

Kurt is one of the Kazakh national products. It is produced by
fermenting cow's, sheep's, or goat's milk with pure lactic acid
streptococcal cultures, separating the whey from the curd, and then
drying it.

Normalized milk with a fat mass fraction of 0.6\% is pasteurized at a
temperature of 80-85°C for 10-20 minutes and cooled to 32-34°C, after
which 5\% starter culture is added and the acidity during fermentation
is brought to 75-76°T. Then the curd is heated to 38-42°C and held for
20-30 minutes to accelerate whey separation, the whey is removed, and
the curd with a moisture content of 76-80\% is pressed for 3-5 hours. If
salted kurt is prepared, the curd is salted before drying. Then the kurt
is shaped into pieces weighing 20-60 g and dried in special drying
chambers at a temperature of 35-40°C. In the finished product, the mass
fraction of fat should be at least 12\%, moisture content 17\%, salt
content not exceeding 2.5\%, and acidity not exceeding 400°T. The shelf
life of fatty kurt does not exceed 3 months, while that of non-fat kurt
does not exceed 9 months.

The national product kurt is divided into three types:\\
- salted, dried curd, to which salt is added and then shaped into balls
or cylindrical forms and dried;

- boiled, dried curd, which is boiled for 2-3 hours and then shaped into
balls or cylindrical forms and dried;

- boiled paste-like mashed kurt intended for adding to soups {[}4{]}.

Rosehip (\emph{Rosa canina} \emph{L.)} grows in all regions of the
world, including Europe, Africa, Central and Western Asia, and Russia,
due to its non-selective nature regarding climate and soil requirements.
In Kazakhstan, 25 species of rosehip are found, four of which are
endemic. Rosehip occurs in the Eastern Kazakh Uplands, Karatau, and
Western Tien Shan regions. They are perennial plants belonging to the
Rosaceae family. \emph{Rosa} has many species and varieties. Their
fruits, in addition to ascorbic acid, contain vitamins P, B1, B2, A, K,
and E, sugars, tannins, pectins, organic acids, flavonoids, pigments,
and salts of iron, manganese, phosphorus, magnesium, and calcium
{[}5{]}.

Ascorbic acid in rosehip exists in reduced and oxidized forms. After
entering the body, it participates in enzymatic processes, stimulates
metabolism, increases resistance to infections, and enhances work
capacity. Another medicinal property of rosehip is its effect on the
blood coagulation system {[}6{]}.

Vitamin C, or ascorbic acid, is one of the vitamins that performs
important functions in various chemical reactions of cellular
metabolism. A deficiency of vitamin C can lead to scurvy.

In addition, vitamin C deficiency may lead to infections, obesity,
cardiovascular diseases, diabetes mellitus, bone diseases, and skin
diseases {[}7{]}.

{\bfseries Materials and methods.} The research objects were camel milk,
cow's milk, dried and powdered rosehip, and the finished product kurt.
Experimental research work was carried out in the workshop
``Experimental production for dairy processing''.

Camel milk (Kyzylorda), cow's milk, and kurt product were selected as
research objects. Each milk sample (200 ml) was aseptically collected
into sterile plastic bottles with screw caps and immediately delivered
to the laboratory in a refrigerated bag at a temperature of 4°C.

For the purpose of vitamin C enrichment, plant raw material rosehip was
selected. This is because the daily vitamin C intake for adult's ranges
from 65 to 139 mg per day.

Research methods: determination of physicochemical parameters of milk,
methods for determining moisture content and dry matter (GOST 17626-81),
methods for determining acidity (GOST 3626-73), determination of the
mass fraction of vitamin C (GOST 30627.2-98), as well as calculation of
the energy value of the kurt product. In the laboratory of S. Seifullin
Kazakh Agrotechnical Research University, the moisture content and
acidity of the finished product were determined. The vitamin C content
in the finished product was studied in the laboratory of Almaty
Technological University {[}8-10{]}.

A tasting evaluation of the finished product was conducted. Ten tasters
were invited for the tasting, and the organoleptic characteristics were
evaluated. Six samples of the finished product were taken for the
tasting.

{\bfseries Results and discussion.} During the research, for the purpose of
enrichment with vitamin C, formulations of the national fermented dairy
product in six different variants (Table 1) and a technological scheme
(Figure 1) were developed.
\end{multicols}

\fig[0.63\textwidth]{p2/image17}[Fig.1 - Flow diagram of the production technology of kurt with the addition of dried powdered rosehip]

\tcap{Table 1 - Formulation of kurt with the addition of dried powdered rosehip (100 g)}
\begin{longtblr}[
  label = none,
  entry = none,
]{
  width = \linewidth,
  colspec = {Q[160]Q[129]Q[129]Q[129]Q[129]Q[129]Q[129]},
  cells = {c},
  cells = {font = \small},
  row{1} = {font=\bfseries\small},
  hlines,
  vlines,
}
{\textbf{Raw~material}\\\textbf{name}} & Sample~No.1 & Sample~No.2 & Sample~No.3 & Sample~No.4 & Sample~No.5 & Sample~No.6\\
{Curd~made~from\\cow’s~milk} & 98 & 95 & 93 & – & – & –\\
{Curd~made~from\\~camel~milk} & – & – & – & – & – & –\\
{Dried~powdered\\~rosehip} & – & 3 & 5 & 98 & 95 & 93\\
Salt & 2 & 2 & 2 & 2 & 3 & 5\\
Total & 100 & 100 & 100 & 100 & 100 & 100
\end{longtblr}

\begin{multicols}{2}
According to the organoleptic indicators, the best sample was selected
depending on the level of vitamin C enrichment.

Figure 1 shows the technology of the national kurt product enriched with
plant-based raw material rosehip.

Kurt should be stored at a temperature not exceeding 15°C and at a
relative air humidity not exceeding 75\%; fatty kurt should be stored
for 1 month, and skimmed (defatted) kurt for 3 months. The
physicochemical parameters of milk are presented in Table 2
\end{multicols}

\tcap{Table 2 - Physicochemical parameters of milk}
\begin{longtblr}[
  label = none,
  entry = none,
]{
  cells = {c},
  cells = {font = \small},
  row{1} = {font=\bfseries\small},
  hlines,
  vlines,
}
Parameters & Control~sample~(cow’s~milk) & Sample~I (skimmed~camel~milk)\\
Acidity, °T & 18 & 19\\
Fat, \% & 3,01 & 1,46\\
Dry~matter, \% & 11,86 & 8,95\\
Solids-not-fat~(SNF), \% & 10,41 & 8,85\\
Protein, \% & 3,33 & 3,30\\
Density,~kg/m³ & 1031,28 & 1033,18
\end{longtblr}

\begin{multicols}{2}
The evaluation of the organoleptic properties of kurt with the addition
of dried powdered rosehip was carried out by assessing aroma, color,
consistency, appearance, and taste. Table 3 presents the results of the
tasting of the finished product, i.e., its organoleptic characteristics.
\end{multicols}

\tcap{Table 3 - Organoleptic characteristics of kurt with the addition of dried powdered rosehip}
\begin{longtblr}[
  label = none,
  entry = none,
]{
  width = \linewidth,
  colspec = {Q[175]Q[360]Q[315]Q[88]},
  cells = {font = \small},
  row{1} = {c,font=\bfseries\small},
  column{4} = {c,font=\bfseries\small},
  hlines,
  vlines,
}
Samples & Appearance and consistency & Taste and aroma & Color\\
{Control sample\\ (without additives)} & {Weight 2-60g, spherical.\\Irregular shapes, with grooves and rounded edges allowed. Hard and dry} & Acidic and moderate & Milky white\\
{Control sample\\(3\% rosehip)} & {Weight 2-60g, spherical.\\Irregular shapes, with grooves and rounded edges allowed. Hard and dry} & Acidic and moderate, slightly salty & Yellowish\\
{Control sample\\ (5\% rosehip)} & {Weight 2-60g, spherical.\\Irregular shapes, with grooves and rounded edges allowed. Hard and dry} & {Acidic and moderate, slightly salty, rosehip taste not\\noticeable} & Light brown\\
{Kurt made from camel milk\\(without additives)} & {Weight 2-60g, spherical.\\Irregular shapes, with grooves and rounded edges allowed. Hard and dry} & Acidic and moderate, slightly salty, rosehip taste weak & White\\
{Kurt made from camel milk\\(3\% rosehip)} & {Weight 2-60g, spherical.\\Irregular shapes, with grooves and rounded edges allowed. Hard and dry} & Acidic and moderate, slightly salty & Yellowish\\
{Kurt made from camel milk\\(5\% dried rosehip)} & {Weight 2-60g, spherical.\\Irregular shapes, with grooves and rounded edges allowed. Hard and dry} & Acidic and moderate,  slightly salty, rosehip taste not noticeable & Light brown
\end{longtblr}

\begin{multicols}{2}
The organoleptic properties were evaluated using a 5-point scale. In the
overall score assessment, the kurt containing 5\% dried powdered rosehip
received the highest scores (4.9 and 4.89) (Figure 2). The control
sample was rated 4,9 and 4,5 points for taste and aroma, while the kurt
with 5\% dried powdered rosehip received 4,7 and 4,7 points,
respectively.

Based on organoleptic properties and vitamin C content, six samples were
selected.

In 100 g of kurt with powdered rosehip, the fat content is 12,1 g,
protein 53,2 g, carbohydrates 17,4 g, and its energy value is 391,3 kcal
or 1637,2 kJ.
\end{multicols}

\fig[0.73\textwidth]{p2/image18}[Fig.2 - Organoleptic score evaluation of kurt with the addition of dried powdered rosehip]

\tcap{Table 4 - Physicochemical parameters of kurt}
\begin{longtblr}[
  label = none,
  entry = none,
]{
  width = \linewidth,
  colspec = {Q[200]Q[265]Q[242]Q[279]},
  cells = {c},
  cells = {font = \small},
  row{1} = {font=\bfseries\small},
  hlines,
  vlines,
}
Parameter & Control sample (without additives) & Control sample (5\% rosehip) & Kurt from camel milk (5\% rosehip)\\
Acidity, °T & 320 & 330 & 350\\
Moisture content, \% & 15 & 17 & 17
\end{longtblr}

\begin{multicols}{2}
In Table 4, the physicochemical parameters of the finished product were
determined according to GOST 17626-81 and GOST 3626-73.

The content of vitamin C in 100 g of the finished product was determined
using the colorimetric method according to GOST 30627.2-98 (Table 5).
\end{multicols}

\tcap{Table 5 - Vitamin C content in 100 g of products}
\begin{longtblr}[
  label = none,
  entry = none,
]{
  cells = {c},
  cells = {font = \small},
  row{1} = {font=\bfseries\small},
  hlines,
  vlines,
}
Sample & Vitamin C content, mg/100 g\\
Control sample (without additives) & 1,13 ± 0,03\\
Control sample (5\% rosehip) & 8,125 ± 0,19\\
Kurt from camel milk (5\% rosehip) & 11,754 ± 0,24
\end{longtblr}

\begin{multicols}{2}
{\bfseries Conclusion.} Based on the results of the conducted research, the
possibility of introducing dried powdered rosehip into production to
preserve the quality of the finished product, enhance its nutritional
value, and enrich it with vitamin C was studied. The technology for
producing kurt with dried and powdered rosehip was developed. The
physicochemical parameters of the kurt product were determined according
to GOST 17626-81 and GOST 3626-73. The organoleptic characteristics,
nutritional, and energy values of the finished product were evaluated.
Based on the research results, the proposed national kurt product
enriched with vitamin C can be considered a natural product recommended
for general consumption.
\end{multicols}

\begin{center}
{\bfseries Литература}
\end{center}

\begin{refs}
1. Baiterek. Concept of development of agro-industrial complex of the
Republic of Kazakhstan for 2021-2030
years.URL:https://baiterek.gov.kz/en/programs/concept-of-development-of-agro-industrial-complex-of-the-republic-of-kazakhstan-for-2021-2030-years.
- Дата обращения: 18.12.2025.

2. Ali A.H., Abu-Jdayil B., Bamigbade G., Kamal-Eldin A., Hamed F.,
Huppertz T., Liu Sh-Q., Ayyash M. Properties of Low-Fat Cheddar Cheese
Prepared from Bovine-Camel Milk Blends: Chemical Composition,
Microstructure, Rheology and Volatile Compounds// Dairy Science. -2024.
-Vol.107(5). - P.2706-2720. DOI 10.3168/jds.2023-23795.

3. Ali A.H., Li S., Liu Sh-Q., Gan R-Y, Li H-B, Kamal-Eldin A., Ayyash
M. Invited Review: Camel Milk and Gut Health: Understanding
Digestibility and the Impact on Gut Microbiota//Dairy Science. - 2024.
-Vol.107(5). - P.2573-2585. DOI 10.3168/jds.2023-23995.

4. Итбалакова А.Б., Ребезев М.Б.Разработка технологии нового
кисломолочного продукта с использованием микробиологического
консорциума//Вестник АТУ. -2025. -T.147(1). -
С.106-114.~\href{https://doi.org/10.48184/2304-568X-2025-1-106-114}{DOI
10.48184/2304-568X-2025-1-106-114}.

5. Алимарданова М.К., Петченко В.И., Елемесова А.А. Совершенствование
технологии кисломолочных продуктов функционального назначения//Продукты
питания: Производство, безопасность, качество. Международной
научно-практической конференции. Башкирский ГАУ. -2019. -С.59-62.

6. Мустафаева А.К., Калемшарив Б, Жандауова А. Түйе сүтінен ұлттық
сүтқышқылды өнімін өндіру технологиясын әзірлеу//Механика және
технологиялар Ғылыми журнал. -2024. - №4 (86). -Б.9-16. DOI
\href{https://doi.org/10.55956/QKOR8581}{10.55956/QKOR8581}.

7. Sahingil D., Hayaloglu A.A. Enrichment of antioxidant activity,
phenolic compounds, volatile composition and sensory properties of
yogurt with rosehip (Rosa canina L.) fortification // International
Journal of Gastronomy and Food Science. -2022.- Vol.28:100514. DOI
\href{https://doi.org/10.1016/j.ijgfs.2022.100514}{10.1016/j.ijgfs.2022.100514}.

8. ГОСТ 17626-81. Казеин технический. Технические
условия/Стандартинформ. -2008. -С.20.

9. ГОСТ 3626-73. Молоко и молочные продукты. Методы определения влаги и
сухого вещества/Стандартинформ. -2009. -С.11.

10. ГОСТ 30627.2-98. Продукты молочные для детского питания. Методы
измерений массовой доли витамина С (аскорбиновой кислоты)/
Межгосударственный совет по стандартизации, метрологии и сертификации.
Минск. -2009. -С.8.
\end{refs}

\begin{center}
{\bfseries References}
\end{center}

\begin{refs}
1. Baiterek. Concept of development of agro-industrial complex of the
Republic of Kazakhstan for 2021-2030
years.URL: \href{https://baiterek.gov.kz/en/programs/concept-of-development-of-agro-industrial-complex-of-the-republic-of-kazakhstan-for-2021-2030-years}{https://baiterek.gov.kz}.
- Date of address: 18.12.2025.

2. Ali A.H., Abu-Jdayil B., Bamigbade G., Kamal-Eldin A., Hamed F.,
Huppertz T., Liu Sh-Q., Ayyash M. Properties of Low-Fat Cheddar Cheese
Prepared from Bovine-Camel Milk Blends: Chemical Composition,
Microstructure, Rheology and Volatile Compounds// Dairy Science. -2024.
-Vol.107(5). - P.2706-2720. DOI 10.3168/jds.2023-23795.

3. Ali A.H., Li S., Liu Sh-Q., Gan R-Y, Li H-B, Kamal-Eldin A., Ayyash
M. Invited Review: Camel Milk and Gut Health: Understanding
Digestibility and the Impact on Gut Microbiota//Dairy Science. - 2024.
-Vol.107(5). - P.2573-2585. DOI 10.3168/jds.2023-23995.

4. Itbalakova A.B., Rebezev M.B. Razrabotka tehnologii novogo
kislomolochnogo produkta s ispol' zovaniem
mikrobiologicheskogo konsorciuma//Vestnik ATU. -2025. -T.147(1). -
S.106-114. DOI 10.48184/2304-568X-2025-1-106-114. {[}in Russian{]}

5. Alimardanova M.K., Petchenko V.I., Elemesova A.A. Sovershenstvovanie
tehnologii kislomolochnyh produktov funkcional' nogo
naznachenija//Produkty pitanija: Proizvodstvo,
bezopasnost', kachestvo. Mezhdunarodnoj
nauchno-prakticheskoj konferencii. Bashkirskij GAU. -2019. -S.59-62.
{[}in Russian{]}

6. Mustafaeva A.K., Kalemshariv B, Zhandauova A. Tүje sүtіnen ұlttyқ
sүtқyshқyldy өnіmіn өndіru tehnologijasyn әzіrleu//Mehanika zhәne
tehnologijalar Ғylymi zhurnal. -2024. - №4 (86). -B.9-16. DOI
10.55956/QKOR858. {[}in Russian{]}

7. Sahingil D., Hayaloglu A.A. Enrichment of antioxidant activity,
phenolic compounds, volatile composition and sensory properties of
yogurt with rosehip (Rosa canina L.) fortification // International
Journal of Gastronomy and Food Science. -2022.- Vol.28:100514. DOI
\href{https://doi.org/10.1016/j.ijgfs.2022.100514}{10.1016/j.ijgfs.2022.100514}.

8. GOST 17626-81. Kazein tehnicheskij. Tehnicheskie
uslovija/Standartinform. -2008. -S.20. {[}in Russian{]}

9. GOST 3626-73. Moloko i molochnye produkty. Metody opredelenija vlagi
i suhogo veshhestva/Standartinform. -2009. -S.11. {[}in Russian{]}

10. GOST 30627.2-98. Produkty molochnye dlja detskogo pitanija. Metody
izmerenij massovoj doli vitamina S (askorbinovoj kisloty)/
Mezhgosudarstvennyj sovet po standartizacii, metrologii i sertifikacii.
Minsk. -2009. -S.8. {[}in Russian{]}
\end{refs}

\begin{info}
\hspace{1em}\emph{{\bfseries Information about the authors}}

Kalemshariv B.-Lecturer of the Department of Food Technology and
Processing Products, Saken Seifullin Kazakh Agrotechnical Research
University, Astana, Kazakhstan, е-mail: begjan.ae@gmail.com;

Askar R. -student, Saken Seifullin Kazakh Agrotechnical Research
University, Astana, Kazakhstan, e-mail: raushan.askar@internet.ru;

Mustafayeva A.K.-Candidate of Technical Sciences, Senior lecturer,
S.Seifullin Kazakh Agrotechnical Research University, Astana,
Kazakhstan; е-mail: ayaulym.mustafa@mail.ru;

Mursalykova M.- PhD doctor, Saken Seifullin Kazakh Agrotechnical
Research University, Astana, Kazakhstan, е-mail: maigul\_85@mail.ru;

Shulenova A. M.-Doctoral student, Saken Seifullin Kazakh Agrotechnical
Research University, Astana, Kazakhstan, e-mail: shulenovaa@mail.ru.

\hspace{1em}\emph{{\bfseries Сведения об авторах}}

Калемшарив Б.-преподаватель, Казахский агротехнический
исследовательский университет им. Сакена Сейфуллина, Астана,
Казахстан, е-mail: begjan.ae@gmail.com;

Аскар Р.-студентка, Казахский агротехнический исследовательский
университет им. Сакена Сейфуллина»,. Астана, Казахстан, email:
raushan.askar@internet.ru;

Мустафаева А.К. -кандидат технических наук, ассоциированный профессор,
Казахский агротехнический исследовательский университет им.
С.Сейфуллина, Астана, Казахстан, е-mail: ayaulym.mustafa@mail.ru;

Мурсалыкова М.Т.-доктор PhD,Казахский агротехнический
исследовательский университет им. Сакена Сейфуллина, Астана,
Казахстан, е-mail:maigul\_85@mail.ru;

Шуленова А.М.-докторант, Казахский агротехнический исследовательский
университет им. Сакена Сейфуллина, Астана, Казахстан,е-mail:
shulenovaa@mail.ru.
\end{info}
